\section{Methods}

\subsection{Study Type}

This is an exploratory evidence map and research-program manuscript. It is not
a systematic review, clinical trial, case report, or clinical guideline. The
purpose is to make the research question, source trail, candidate gates, and
next experiments inspectable.

\subsection{Source Intake}

The repository uses retrieval-led evidence collection. Priority sources include
clinical guidelines, PubMed-indexed literature, ClinicalTrials.gov records,
regulatory pages, official chemistry databases, public full-text sources, and
classical Islamic sources when they are used for ethics or hypothesis framing.

Each reusable finding is stored as a dated note under \repoPath{research/}.
Small reproducible source snapshots are stored under \repoPath{data/}. The
thalassemia source registry links claims to the evidence trail.

\subsection{Evidence Classification}

Evidence is classified into five practical labels:

\begin{itemize}[leftmargin=*]
  \item \textbf{Benchmark:} an existing therapy or endpoint used for comparison.
  \item \textbf{Candidate:} a molecule, material, or mechanism that needs
  validation before treatment language is allowed.
  \item \textbf{Gate:} a required decision point such as subtype, genotype,
  immune status, organ risk, safety, access, or clinician review.
  \item \textbf{Boundary:} evidence that prevents overclaiming.
  \item \textbf{Open question:} a high-value unknown that needs records, public
  data, or wet-lab work.
\end{itemize}

\subsection{AI-Assisted Workflow}

AI was used to retrieve, summarize, cross-link, and draft research artifacts.
The repository keeps source snapshots and dated notes so claims can be audited
without trusting the model output alone. AI-generated text is treated as a
drafting and organization aid; scientific and clinical claims remain bound to
source evidence and clinician or domain-expert review.

\subsection{No-Lab Constraint}

Nakafa Lab does not currently operate a wet lab or clinical trial. Work that can
continue without a lab partner includes source-backed literature mapping,
computational candidate rejection, chemistry identity checks, de-identified
case-packet preparation, assay design, and specialist-referral question
generation. Patient-specific treatment decisions are out of scope.
